\section{Introduction}
\label{main:introduction}

Let $h=(h(1),\ldots,h(n))$ be a weakly increasing integer sequence with
$i\le h(i)\le n$. The associated \emph{natural unit interval graph} has
vertex set $[n]=\{1,\ldots,n\}$ and edges
\[
 E(G)=\bigl\{\{i,j\}:i<j\le h(i)\bigr\}.
\]
These are also called naturally labeled Dyck graphs. A proper coloring is a
map $f:[n]\to\N_{>0}$ that has distinct values on adjacent vertices. In the
variable alphabet $X=(x_1,x_2,\ldots)$ put
\begin{equation}
 \chi_G(X;q)=\sum_{f\text{ proper}}
 q^{\operatorname{asc}_G(f)}\prod_{a=1}^n x_{f(a)},
 \qquad
 \operatorname{asc}_G(f)=
 \#\{\{a,b\}\in E(G):a<b,\ f(a)<f(b)\}.
 \label{main:csf}
\end{equation}
Shareshian and Wachs proved that this function is symmetric for this graph
class and conjectured that its coefficients in the elementary basis belong
to $\N[q]$ \cite{SW16}. Here
$e_k(X)=\sum_{i_1<\cdots<i_k}x_{i_1}\cdots x_{i_k}$ and
$e_\lambda=\prod_i e_{\lambda_i}$ for a partition $\lambda$.

The coefficients admit a more precise interpretation in terms of
permutations. Define
\begin{align*}
 D_G^0&=\{\sigma\in S_n:
   \sigma_i>\sigma_{i+1}\Longrightarrow
   \{\sigma_i,\sigma_{i+1}\}\in E(G)\},\\
 \ginv_G(\sigma)&=\#\{(i,j):i<j,\ \sigma_i>\sigma_j,
                    \{\sigma_i,\sigma_j\}\in E(G)\}.
\end{align*}
Thus $D_G^0$ forbids adjacent descents between nonneighbors, and
$\ginv_G$ counts only inversions supported on edges.

\begin{theorem}[Elementary positivity and a permutation witness]
\label{main:theorem}
There is a deterministic terminating procedure which, given a natural
unit interval graph $G$ on $[n]$ and $\sigma\in D_G^0$, returns a partition
$\theta_G(\sigma)\vdash n$ such that
\begin{equation}
 \chi_G(X;q)=\sum_{\sigma\in D_G^0}
       q^{\ginv_G(\sigma)}e_{\theta_G(\sigma)}(X).
 \label{main:witness}
\end{equation}
In particular every elementary coefficient of $\chi_G(X;q)$ belongs to
$\N[q]$. The assertion holds as a formal symmetric-function identity,
uniformly in $n$ and $G$.
\end{theorem}

This proves the elementary-positivity part of the Shareshian--Wachs
conjecture. The theorem makes no assertion of elementary unimodality.
The construction specifies a map on individual inputs: it does not take
the elementary coefficient list as input and redistribute permutations
according to that list. Its running time may nevertheless be large.

For example, let $G$ be the path $1$--$2$--$3$. Its expansion is
\begin{equation}
 \chi_G(X;q)=q\,e_{(2,1)}(X)+(1+q+q^2)e_{(3)}(X).
 \label{main:path-example}
\end{equation}
Write $m_\lambda$ for the sum of the distinct monomials with exponent
multiset $\lambda$. To check this directly, the colorings with a repeated color contribute
$q\,m_{(2,1)}(X)$, while those with three distinct colors contribute
$(1+4q+q^2)m_{(1,1,1)}(X)$. Now use
$e_{(2,1)}=m_{(2,1)}+3m_{(1,1,1)}$ and
$e_{(3)}=m_{(1,1,1)}$.
The domain $D_G^0$ consists of $123,132,213,321$, with graph-inversion
degrees $0,1,1,2$. The theorem assigns one of the two degree-one
permutations to $(2,1)$ and the other three permutations to $(3)$.
Thus the map refines the degree count by an entire elementary partition.
Equation~\eqref{main:path-example} alone does not specify which of
$132$ and $213$ receives which partition; the procedure in
Section~\ref{route:section} makes that individual choice.

The positivity question is difficult because the coloring definition
already has nonnegative coefficients in the monomial variables, whereas
rewriting it in the elementary basis can involve subtraction.  Our proof
separates this issue from the individual assignment.  It first identifies
each relevant elementary coefficient with a positive algebraic object.
It then realizes the algebraic identities by finite bijections between
ordered collections of independent sets. Neither stage uses the sought elementary
coefficient list as input to the permutation rule.

\subsection{History and significance}

The problem grew out of positivity questions for immanants of
Jacobi--Trudi matrices, studied by Haiman and by Stanley and Stembridge
\cite{Hai93,SS93}.  Stanley introduced the chromatic symmetric function
as a symmetric-function refinement of the chromatic polynomial
\cite{Sta95}; it is the specialization of
Equation~\eqref{main:csf} at $q=1$.  In this language, the
Stanley--Stembridge conjecture asks for elementary positivity of the
incomparability graphs of $(3+1)$-free posets
\cite[Conjecture~5.5]{SS93}. The incomparability graph joins two
elements exactly when neither is below the other. A poset is
$(3+1)$-free if it has no induced three-element chain together with
an element incomparable to that chain.
Guay-Paquet's reduction expresses its chromatic symmetric function as
a nonnegative rational combination of those of unit interval orders
\cite[Theorem~5.1]{GP13}. Their naturally labeled incomparability
graphs are the graphs in Theorem~\ref{main:theorem}. Specializing our
theorem at $q=1$ therefore also proves the classical positivity
theorem, previously established by Hikita \cite{Hik24}.

Positive expansions in other bases provided important earlier
structure. Haiman's immanant results imply Schur positivity for
incomparability graphs of $(3+1)$-free posets \cite{Hai93}; Gasharov
gave a positive formula using poset tableaux \cite{Gas96}. Chow's
quasisymmetric expansion related chromatic symmetric functions to
permutation descents \cite{Cho99}.  Shareshian and Wachs introduced the
$q$-refinement in Equation~\eqref{main:csf}, proved its symmetry and
Schur positivity for natural unit interval graphs, and formulated the
elementary-positivity and elementary-unimodality conjecture
\cite{SW12,SW16}.  Their acyclic-orientation formula gives elementary
coefficients after grouping partitions by their number of parts
\cite[Theorem~5.3]{SW16}; Athanasiadis supplied a power-sum expansion
\cite{Ath15}.  These descriptions do not determine a nonnegative
coefficient formula for each elementary partition.

The geometric connection gives another reason to study the
$q$-grading.  Tymoczko constructed permutation actions on equivariant
cohomology \cite{Tym08}, and Shareshian and Wachs related their
chromatic functions conjecturally to the graded representation on
regular semisimple Hessenberg varieties \cite{SW12,SW16}.
Brosnan--Chow and Guay-Paquet proved this representation-theoretic
identity \cite{BC18,GP16}.  That result is distinct from the
elementary-positivity assertion addressed here.  Geometric,
combinatorial and algebraic approaches subsequently established
positivity for substantial families: abelian Hessenberg varieties
\cite{HP19}, classes treated by explicit elementary expansions
\cite{CH19}, melting lollipops \cite{HNY20}, and the $q$-rook
formulas for abelian Dyck paths \cite{CMP23}.  The modular law also
provided a systematic tool for organizing chromatic functions and
their relations \cite{AN21}.

Hikita's proof uses a probabilistic formula \cite{Hik24}. His formula gives numerical
nonnegativity for every positive real $q$, while its rational
$q$-weights do not establish coefficientwise nonnegativity in $\N[q]$;
the refined conjecture remains explicitly distinguished in
\cite[Section~1.3]{Hik24}.  Griffin, Mellit, Romero, Weigl and Wen gave
an independent proof through Macdonald expansions, recovering Hikita's
elementary formula \cite{GMRWW25}. Further developments include
positivity at $q=1$ of the Abreu--Nigro $g$-functions, proved by
Huh, Hwang, Kim, Kim and Oh as a refinement of Hikita's theorem
\cite{HHKKO25},
Chow's combinatorial interpretation of the probabilities \cite{Chow26},
and lower bounds for elementary coefficients studied by Siegl
\cite{Siegl26}.  Theorem~\ref{main:theorem} gives a formal polynomial
formula with nonnegative integer coefficients and an individual
assignment preserving the specified graph-inversion weight.

The representation identity and Theorem~\ref{main:theorem} give the
following degreewise consequence.

\begin{corollary}[Young permutation modules in Hessenberg cohomology]
\label{main:hessenberg}
Let $n\ge1$, let $h$ be as above with graph $G=G(h)$, and let
$S\in M_n(\C)$ be regular semisimple. Write
\[
 \operatorname{Hess}(S,h)=
 \{F_\bullet\in\operatorname{Flag}(\C^n):
     SF_i\subseteq F_{h(i)}\text{ for all }i\}.
\]
For $\lambda=(\lambda_1,\ldots,\lambda_\ell)\vdash n$, let
$M^\lambda=\operatorname{Ind}_{S_\lambda}^{S_n}\mathbf{1}_{\C}$ be the
Young permutation module, where
$S_\lambda=S_{\lambda_1}\times\cdots\times S_{\lambda_\ell}$ is the
standard Young subgroup, and put
\[
 m_{\lambda,d}(h)=
 \#\{\sigma\in D_G^0:\theta_G(\sigma)=\lambda,\ \ginv_G(\sigma)=d\}.
\]
Then for every integer $d\ge0$, under Tymoczko's dot action on ordinary
cohomology there is an abstract $\C[S_n]$-module isomorphism
\[
 H^{2d}(\operatorname{Hess}(S,h);\C)
 \cong\bigoplus_{\lambda\vdash n}(M^\lambda)^{\oplus m_{\lambda,d}(h)}.
\]
\end{corollary}

\begin{proof}
Write $\operatorname{ch}$ for the Frobenius characteristic and $\omega$
for the usual symmetric-function involution, so
$\omega(e_\lambda)=h_\lambda$, the product of complete homogeneous
symmetric functions. The identity of Brosnan--Chow
\cite[Theorem~129]{BC18} and Guay-Paquet \cite{GP16}, followed by
Equation~\eqref{main:witness}, gives
\[
 \operatorname{ch}H^{2d}(\operatorname{Hess}(S,h);\C)
 =[q^d]\,\omega\chi_G(X;q)
 =\sum_{\lambda\vdash n}m_{\lambda,d}(h)h_\lambda.
\]
Since $\operatorname{ch}M^\lambda=h_\lambda$
\cite[proof of Proposition~10]{BC18}, injectivity of the Frobenius
characteristic and complete reducibility of finite-dimensional complex
$S_n$-modules give the isomorphism.
The argument supplies no preferred isomorphism, geometric basis, or
compatibility with cup product.
\end{proof}

There is also a direct precedent for asking for a partition-valued
permutation rule.  Stanley and Stembridge proposed such an assignment
on their compatible permutations \cite[Remark~5.6]{SS93}.  That proposal
concerns the ungraded elementary expansion.  The rule here retains the
full graph-inversion degree and specifies a terminating procedure on
each input.  The finite cancellation used to pass from required
ascents to exact ascents belongs to the bijective-cancellation
tradition represented by the Garsia--Milne involution principle
\cite{GM81}.

Our algebraic reduction likewise has substantial antecedents.
Stanley's graph specialization \cite[Section~2]{Sta98}, chromatic functions in
noncommuting variables \cite{GS01}, and the noncommutative
symmetric-function methods of Fomin--Greene and Blasiak--Fomin
\cite{FG98,BF17} provide a framework for commuting symmetric expressions
inside a noncommutative algebra.  Hwang's noncommutative
$P$-symmetric functions \cite{Hwa24} and the poset algebras of Blasiak,
Eriksson, Pylyavskyy and Siegl \cite{BEPS25} provide direct predecessors
of the graph-alphabet evaluation and its elementary/monomial
coefficient pairing.  We give a short proof in the centered quantum
torus conventions needed for the subsequent argument.

The positive factorization and transport stages use the mathematical
framework of Hall algebras and scattering diagrams.  Ordered slope
factorization and polynomial shuffle products were developed in
\cite{Rei03,KS11}, with semistable quotient and integrality results in
\cite{Ef12,FR18,DMe20}.  The connection from Hall factorization to
wall transformations appears in Reineke's Poisson automorphisms and
Bridgeland's stability scattering diagrams \cite{Rei10,Bri17}.
Scattering constructions originate in work of Kontsevich--Soibelman
and Gross--Siebert \cite{KS06,GS11,KS14}; theta bases and mutation are
central to the cluster framework of Gross--Hacking--Keel--Kontsevich
\cite{GHKK18}, and Davison--Mandel prove strong quantum positivity
\cite{DM21}.  We prove the particular reordering, completion and
transport statements needed here with their explicit gradings and
finite coefficient bounds.  The graph-specific step is then an
identification: general positivity of theta sections does not itself
show that a graph-alphabet monomial symmetric function is one such
section.  The triangular chart below proves that identity before the
packet construction extracts the individual witness.

\subsection{Proof structure}

The first reduction replaces colorings by coefficients in an algebra
whose monomials record vertex multiplicities. Write $q=v^2$ and let
$X_m$ denote the basis monomial indexed by a vertex-multiplicity
vector $m$.  Multiplication of these monomials is
twisted by an alternating form recording ordered edges. For a larger unit
interval graph $\Gamma$ containing $G$, let $E_k$ be the sum of the
monomials of independent $k$-sets. We prove that the $E_k$ commute, so
they may be evaluated as elementary symmetric functions $e_k(Y)$ of a
formal alphabet $Y$. If $u$ is the sum of the original vertex exponents
and $M=|E(G)|$, the elementary coefficient becomes
\begin{equation}
 [e_\lambda(X)]\chi_G(X;v^2)
       =v^M[X_u]m_\lambda(Y).
 \label{main:reduction}
\end{equation}
Here $m_\lambda$ is the monomial symmetric function. The identity is
proved in Proposition~\ref{graph:kernel}.

The central step is to identify $m_\lambda(Y)$ as a chamber value of a
single positive section of a system of wall transformations. A wall is
a portion of a hyperplane in the real dual of the exponent lattice,
equipped with conjugation by a formal series. A section is a family of
chamber values related by these transformations. We construct sections $\Theta_m$,
indexed by lattice exponents, by prescribing their coefficients in test
directions determined by the alternating form. The auxiliary graph is
arranged in a triangle with distinguished independent vertices
$d_1,\ldots,d_D$, called anchors. For a partition $\lambda$ of length
at most $D$, padded by zeroes to length $D$, the essential identification in the all-negative chamber is
\begin{equation}
 m_\lambda(Y)=\Theta_{\lambda_1d_1+\cdots+\lambda_Dd_D}.
 \label{main:identification}
\end{equation}
Repeated interval rotations establish this identity. A support argument
using positive powers of $E_1$ is essential: stationarity under the
rotations alone does not force an exponent to involve anchors only.

The proof has three main components.  First, positive quantum factors
remain positive when their order is reversed.  We realize the new
multiplicities as dimensions of explicit graded vector spaces, and
prove the coefficient bounds needed for all infinite series.
Second, the interval rotations identify a family of graph symmetric
functions with single positive theta sections. Together with
Equation~\eqref{main:reduction}, this already proves coefficientwise
elementary positivity. Third, we extract the individual assignment.
After reversing an input permutation, we group vertices into packets,
each an independent set written in increasing order. A local packet
bijection transports them to subsets of the independent anchors and
records weighted jumps. The total jump histories are independent of
packet sizes. Finite matching cancellation refines required nonedge
ascents of input words and required strict ascents of anchor words to
exact sets of such positions. The empty set on the anchor side leaves
a weakly decreasing word, whose multiplicities give the output partition. The finite
matching argument supplies termination without consulting the desired
final shape counts.

Section~\ref{graph:section} proves the coloring reduction and fixes
permutation conventions. Section~\ref{alg:section} proves the
reordering result. Section~\ref{wall:section} constructs the wall
sections and proves their positivity and mutation rules.
Section~\ref{tri:section} establishes
Equation~\eqref{main:identification}, and
Section~\ref{route:section} constructs the individual permutation map.
Section~\ref{main:completion} returns to the original graph and proves
Theorem~\ref{main:theorem}, including the reversal needed for the
convention in $D_G^0$.
